% O número de permutações de n elementos.
% Compilação: latexmk  (gera O-Numero-de-Permutacoes.pdf)
\documentclass[a4paper,12pt]{article}

\usepackage{fontspec}
\usepackage[brazilian]{babel}
\usepackage[hmargin=2.7cm,vmargin=2.5cm]{geometry}
\usepackage{amsmath,amssymb,mathtools,amsthm}
\usepackage{microtype}
\usepackage{csquotes}
\usepackage[hidelinks,pdftitle={O número de permutações},
  pdfauthor={Cauê Takahashi de Oliveira}]{hyperref}

\newtheorem*{proposicao}{Proposição}

\setlength{\parindent}{0pt}
\setlength{\parskip}{0.6em}

\title{O número de permutações de $n$ elementos}
\author{Cauê Takahashi de Oliveira}
\date{Novembro de 2026}

\begin{document}
\maketitle
\thispagestyle{empty}

\begin{proposicao}
  Seja $A_n = \{a_1, a_2, \dots, a_n\}$ um conjunto com $n \ge 1$ elementos e
  $S_n$ o conjunto das bijeções de $A_n$ em $A_n$. Então $|S_n| = n!$.
\end{proposicao}

\textit{Demonstração.}
Seja $\sigma^{(n)}\colon A_n \to A_n$, $A_n = \{a_1, a_2, \dots, a_n\}$ e
$\sigma^{(n)}$ bijetora;
\[
  S_n \coloneqq \bigl\{\sigma^{(n)} \bigm| \sigma^{(n)} \text{ bijetora}\bigr\}.
\]
Seja $S_n^{j,i} \subset S_n$,
\[
  S_n^{j,i} \coloneqq \bigl\{\sigma \in S_n \bigm| \sigma(a_j) = a_i\bigr\},
  \qquad i, j \in \{1, \dots, n\} \text{ fixos}.
\]
É óbvio que\footnote{Além disso, para $j$ fixo, os conjuntos $S_n^{j,i}$ são
  disjuntos dois a dois, pois $\sigma(a_j)$ assume um único valor; é isso que
  permite somar as cardinalidades mais adiante.}
\[
  \bigcup_{i=1}^{n} S_n^{j,i} = S_n, \qquad \forall j \in \{1, \dots, n\}.
\]
Por simetria $\bigl|S_n^{j,i}\bigr| = \bigl|S_n^{j',i'}\bigr|$,
$\forall\, i, j, i', j' \in \{1, \dots, n\}$.

Seja então\footnote{Aqui $A_{n+1} = A_n \cup \{a_{n+1}\}$. Vendo cada função
  como o conjunto dos pares $(a, \sigma(a))$, a condição
  $\sigma^{(n+1)}|_{A_n} = \sigma^{(n)}$ é $\sigma^{(n)} \subset \sigma^{(n+1)}$,
  ou seja, $\sigma^{(n)} \cap \sigma^{(n+1)} = \sigma^{(n)}$: $\sigma^{(n+1)}$
  estende $\sigma^{(n)}$.}
\[
  \lambda_{n,n+1} \coloneqq \bigl\{(\sigma^{(n)}, \sigma^{(n+1)}) \in S_n \times S_{n+1}
    \bigm| \sigma^{(n+1)}|_{A_n} = \sigma^{(n)}\bigr\}.
\]
Note que $(\sigma^{(n)}, \sigma^{(n+1)}) \in \lambda_{n,n+1} \Rightarrow
\sigma^{(n+1)}(a_{n+1}) = a_{n+1}$. Fixado $\sigma^{(n+1)}(a_{n+1}) = a_{n+1}$,
\begin{align*}
  &\Longrightarrow\ \forall\, \sigma^{(n)} \in S_n,\ \exists!\, \sigma^{(n+1)} \in S_{n+1}
    \bigm| \sigma^{(n+1)}(a_k) = \sigma^{(n)}(a_k),\ \ \forall k \in \{1, \dots, n\}.
  \intertext{Reciprocamente, se $\sigma^{(n+1)}(a_{n+1}) = a_{n+1}$, a restrição
    $\sigma^{(n+1)}|_{A_n}$ é bijetora em $A_n$ (pois $\sigma^{(n+1)}$ é injetora e
    $a_{n+1}$ já é imagem de $a_{n+1}$); logo
    $(\sigma^{(n+1)}|_{A_n}, \sigma^{(n+1)}) \in \lambda_{n,n+1}$.}
  &\Longrightarrow\ \bigl|\lambda_{n,n+1}\bigr| = |S_n|
    = \bigl|S_{n+1}^{n+1,n+1}\bigr| = \bigl|S_{n+1}^{n+1,i}\bigr|,
    \quad \forall i \in \{1, \dots, n+1\}\\
  &\Longrightarrow\ |S_{n+1}| = \Bigl|\bigcup_{i=1}^{n+1} S_{n+1}^{n+1,i}\Bigr|
    = \sum_{i=1}^{n+1} \bigl|S_{n+1}^{n+1,i}\bigr| = (n+1)\cdot|S_n|.
\end{align*}
Como $|S_1| = 1$ e $|S_{n+1}| = (n+1)\cdot|S_n|$, por
PIF\footnote{Princípio da Indução Finita. A base vale porque a única função de
  $\{a_1\}$ em $\{a_1\}$ é a identidade, que é bijetora.}
\[
  |S_1| = 1 \iff |S_2| = 2 \iff \cdots \iff |S_n| = n! \qquad \text{C.Q.D.}
\]

\end{document}
